\documentclass{article}
\usepackage{amsmath}
\begin{document}
\title{A controlled camera displacement restores a hidden cue only inside a calibrated aperture}
\maketitle

\paragraph{Model and comparison.}
Retain the rendered far-site target at depth $Z=37$ or $38$,
half-width $.28$, behind the nearer panel at depth 8 with half-width
$.20$. Both lie on the same horizontal centerline, and the selected
vertical row stays inside their vertical extents. The camera is
stationary during capture. A dedicated station actuator may shift it
right from rest, within a separate budget of $.2$ time units before
the formation's stationary interval. This is modeled station hardware
independent of the mobile member's longitudinal control. This is not
free robot motion or an existing CoSTAR slider. During this single
capture set, there are no earlier preview images or alternate sensors.
Compare exact same-output consumers to an explicitly authorized
producer read at a named ray, retaining the old depth emission and old
radio report.

\paragraph{Sampled aperture and physical lower bound.}
Let $a$ be the pinned horizontal half-projection scale, $W=320$, and
$s_i=a(2i+1-W)/W$ the ray slope in column $i$. At station offset $b$,
the ray hits the target plane horizontally at $b+Zs_i$ and crosses
the occluder plane at $b+8s_i$. For nonnegative shifts up to $.4$ any
ray hitting the target cannot bypass the left side of the panel:
its lateral coordinate at depth 8 is at least $-8(.28)/37$.
Hence the open aperture for a ray supplying a positive qualified
read in both worlds is exactly
\[
  \max\{0,.20-8s_i,-.28-37s_i,-.28-38s_i\}
   < b <
  \min\{.4,.28-37s_i,.28-38s_i\}
\]
in the interior of the modeled domain (the endpoint $.4$ itself is
not operationally used). Geometric grazing boundaries are treated
as nonqualifying. Test the finite ray set directly. The first usable
ray is column 160, with interval approximately
$(.185566202,.211439462)$. The continuous-ray lower threshold for
a common ray would be $.1786667$: choose the largest ray slope
that can still hit the farther target, and compare its crossing at
depth 8 to $.20$. The sampled lower edge is therefore not just the
continuous geometry threshold.

For any column or world, the earliest positive target return is
still above $.18556$ in this pin. In particular the world frames
fully agree where $|b|$ is below the matching first edge (negative
side by symmetry of the target and panel). This is an actual
two-world obstruction to any decoder of the rendered frame there,
not merely to the nominated single-ray algorithm. The independent
old source patch is identical across both worlds and does not
distinguish them.

\paragraph{Add the charged motion and fixed encoder uncertainty.}
With station actuator acceleration bound $A$ and speed bound $V$,
the maximum displacement with rest at both ends over $T$ is the
integral of $\min(At,V,A(T-t))$. Thus at $T=.2,A=20$, speed 2
permits $.2$, whereas speed 1 permits only $.15$. The slower
actuator cannot reach any open view here even with exact pose.
The faster actuator can reach the midpoint of column 160's interval,
at approximately $.198502832$ (rest-to-rest binding is attainable
by a slower triangular stroke). With fixed signed encoder bias of
magnitude at most $.010$, its whole aperture stays valid. The
sharp error supremum for this allowed fixed-ray read is half the
width of that interval, approximately $.012936630$. The other
admissible rays have lower edges above $.2$. State error here is a
constant physical calibration offset, not within-dwell camera
motion or error in the raster depth value. At bias bound $.015$,
every reachable command admits bias $-.015$, putting the camera
at most $.185$. Both worlds then remain hidden, so another
single-frame decoder cannot restore robust service. A proof for
continuous preview, alternate hardware or changed time budget is
outside this comparison.

\paragraph{Selected actual source and timing effects.}
Check 22 real software-rendered frames, including the attempted
low-speed endpoint and both ends of the faster bias interval, with
1,689,600 independent ray comparisons. At the useful position the
same scoped raw read restores the target without suppressing the
old image source emission; the stock point output still clips it.
Feed those rendered pixels into the unmodified selected native
callback/body code and the same broker calendar. The 60 cases
retain the old native radio receipt; the 24 timely valid cues
include only the visible bias settings and eligible budgets.
Neither added airtime nor sharper source interpretation alone
repairs an unreachable camera position. An equally informed
direct controller and slot checker ties; this costs actual
modeled displacement and source read access, not only a new label.
The selected optics is authored, not an actual field camera or an
independently selected cave mission. No full stack acceptance.
\end{document}
